\begin{definition}[The two-gather round]
  \label{def:aba-round-over-gather-specifications}
  \lean{PLTS.ABA.GBCA.candidate}\lean{PLTS.ABA.GBCA.gradeOf}\lean{PLTS.ABA.GBCA.boundOfCore}\lean{PLTS.ABA.GBCA.ByAFW.RoundStateOverGatherSpecifications}\lean{PLTS.ABA.GBCA.ByAFW.AlgorithmOverGatherSpecifications}\lean{PLTS.ABA.GBCA.ByAFW.roundOverGatherSpecifications}\lean{PLTS.ABA.GBCA.ByAFW.roundOverGatherSpecifications_step_iff_algorithm}\leanok
  \uses{def:aba-round-composition, def:aba-gbca-network, def:aba-gather-spec, def:aba-gbca-spec, def:aba-labels}
  The graded agreement of~\cite{AFW25} --- its Algorithm~4 in its case $R \ge 2$, at $R = 2$, with the grade computed from the second gather's counts in place of the approximate-agreement subroutine of lines~7 and~8 (D24) --- over two gather specification instances of Definition~\ref{def:aba-gather-spec}, each the binding form (D33), the first at payload type $\mathbb{B}$ and the second at $\mathbb{B} \cup \{\bot\}$.
  The round is the composition of Definition~\ref{def:aba-round-composition}: its state is the $n$ programs' variables beside the round's network, in parallel with the two gather states, and a program's variables hold the process's input, its candidate, whether it has called the second gather, its graded outcome and its return flag.
  For a first-gather output $g$, the candidate is
  \[
    \mathrm{candidate}(g) \;=\;
      \begin{cases}
        v & \text{if } |{\rm dom}\, g| - f \le \#\{k : g(k) = v\},\\
        \bot & \text{if neither bit qualifies,}
      \end{cases}
  \]
  at most one bit qualifying since the return domains carry $n - f > 2f$ entries; for a second-gather output $h$, the grade is $(v, 2)$ at $|{\rm dom}\, h| - f$ entries of $v$, $(v, 1)$ at $f + 1$, and $(\bot, 0)$ below both.
  The algorithm lists the round's transitions at the specification's labels, one case each: the round's $\mathrm{callG}$ is the caller recording its input beside the first gather's call at a caller holding no input, and the caller's loop beside the first gather's call at one holding an input; each gather's internal transitions are silent transitions; the first gather's return to a process, which records the candidate, the first gather's announcement of its core, which the round's network records, and that process's call of the second gather at the candidate are silent transitions; the second gather's return, which records the grade, and the round's $\mathrm{retG}$, which returns it, are two transitions likewise; a gather's return that the caller's guards refuse is taken on a loop; the announcement fires on $\mathrm{boundG}(r,b)$; and corruption is the simultaneous transform of both gathers (D1).
  A gather specification takes its call label on two transitions, so each of the three transitions reaching a gather call carries the gather's step as a hypothesis, whichever transition it takes.
  The round's network holds the first gather's core it has recorded and whether the round has returned, auxiliary state (D29) that no program reads and no guard consults but the announcement's.
  It records the core from the first gather's announcement, once, and the announcement names $\mathrm{boundOfCore}$ of that core once the round has returned.
  For a payload set $S$, $\mathrm{boundOfCore}(S)$ is the bit carried by at least $|S| - f$ of the entries of $S$, and $1$ where neither bit is.
  At the sizes the core transition of Definition~\ref{def:aba-gather-spec} allows, $|S| \ge n - f > 2f$, at most one bit is carried by that many, so a return of $v$ returns the announced bit, and a round that returns no value announces a bit whose complement no return can have as its value.
  Full construction: \leandecl{PLTS.ABA.GBCA.ByAFW.roundOverGatherSpecifications}; its transitions at the specification's labels are listed by \leandecl{PLTS.ABA.GBCA.ByAFW.AlgorithmOverGatherSpecifications}, a relation on the round's state, and \leandecl{PLTS.ABA.GBCA.ByAFW.roundOverGatherSpecifications_step_iff_algorithm} states that the round's transitions over the family labels the pullback sends to one specification label are exactly the cases of that algorithm at that label.
\end{definition}
